\subsection{Two-Phase Execution System}
We organize execution as a two-phase decomposition. Offline, the system compiles the dataset $\mathcal D$ together with a specification space $\mathcal W$ into reusable structures that support efficient evaluation of admissibility and cost under different $\omega \in \mathcal W$. Online, given a request $(q,\omega,\Omega,k)$ with a concrete $\omega$, the system retrieves a feasible candidate pool under $\omega$ and selects a size-$k$ set by optimizing $\Omega$ while accounting for transformation cost. Figure~\ref{fig:cf_two_stage_pipeline} summarizes the pipeline. Sections~\ref{sec:instance} and~\ref{sec:set} provide domain-specific examples and practical considerations.

\paragraph{Offline (preprocessing for $(\mathcal D,\mathcal W)$).}
Offline, the system takes as input the dataset $\mathcal D$ and a family of admissibility specifications $\omega \in \mathcal W$.  Rather than fixing a single $\omega$, preprocessing prepares structures that remain valid across different instance-level constraints and transformation cost configurations specified at query time. The preprocessing phase produces reusable artifacts:
\begin{itemize}
    \item \textbf{Step 1.a (Admissibility and cost parameterization).} 
    Identify the structural primitives that govern feasibility and transformation cost under $\omega \in \mathcal W$. This includes constraint types, immutable attributes, categorical rules, and cost components that will later instantiate $f_\omega(*,x)$ and $c_\omega(*,x)$ at query time.
    \item \textbf{Step 1.b (Constraint-aware index construction).} Construct an index $\mathcal I$ over $\mathcal D$ to support efficient query under the cost $c_\omega(,x)$ while enforcing the feasibility $f_\omega(,x)$. The index is organized into multiple internal components, each corresponding to a subset of parameters $\omega \in \mathcal W$ that lead to similar query behavior. At query time, the system selects the component whose structure aligns with the cost and constraint by $\omega$.
\end{itemize}
If $\omega$ acts primarily through feature weights that modulate the relative importance of dimensions in $c_\omega$, one may construct metric-aware data structures such as weighted KD-trees~\cite{reif2022scalable}, ball trees~\cite{dolatshah2015ball}, or cover trees~\cite{izbicki2015faster} that support efficient lower-bound pruning under reweighted costs. If instead $\omega$ operates through indicator-type constraints, such as immutability or categorical admissibility rules, one may employ constraint-partitioned indices~\cite{lee2019constraint}, bitmap filtering~\cite{chan1998bitmap}, or inverted indices~\cite{babenko2014inverted} to pre-filter infeasible candidates before cost evaluation. In general, different structural roles of $\omega$ induce different index design choices. 

\paragraph{Online (querying over $(q,\omega,\Omega,k)$).}
In the online stage, directly optimizing Eq.~\eqref{eqn:objective} entails combinatorial search over subsets for a given $(q,\omega,\Omega,k)$, which is computationally expensive and typically relies on heuristics~\cite{qin2009keyword,roy2014fast,yang2016fast}. To reduce complexity and exploit the precompiled index $\mathcal I$, we adopt a coarse-to-fine strategy. Since $\mathcal I$ is organized with respect to $c_\omega(x_q,x)$, it aims to retrieve candidates in nondecreasing cost order, providing a lower bound on the transformation-cost.\footnote{If a query specifies $\omega \notin \mathcal W$, the system cannot leverage the precomputed index $\mathcal I$ for accelerated retrieval; however, counterfactual query remains well-defined and can be performed via a linear scan over $\mathcal D$.}  This enables early pruning before evaluating the set-level objective $\Omega$. The system therefore proceeds as follows.
\begin{itemize}
    \item \textbf{Step 2.a (Cost-aware candidate generation).} Using $\mathcal I$, we construct a candidate pool $\mathcal A_\omega(q)\subseteq \mathcal F_\omega(q)$ by prioritizing low-cost feasible instances under $c_\omega(x_q,x)$. That is,
    \begingroup\small
    \begin{equation}
    \mathcal A_\omega(q) 
    \in 
    \arg\min_{\substack{\mathcal A_\omega \subseteq \mathcal F_\omega(q)\\ |\mathcal A_\omega|=m}}
    \sum_{x\in \mathcal A_\omega} c_\omega(x_q,x).
    \end{equation}
    \endgroup
    Here $m \ge k$ controls the pool size, ensuring that $\mathcal A_\omega(q)$ contains admissible and sufficiently low-effort transformations.
    \item \textbf{Step 2.b (Set selection under $\Omega$).} From the candidate pool $\mathcal A_\omega(q)$, select the final response set $S^\star$ by optimizing the set-level objective $\Omega$, which captures collective criteria such as diversity~\cite{laugel2023achieving}, coverage~\cite{wen2025enhancing}, robustness~\cite{guyomard2023generating}, or fairness~\cite{kusner2017counterfactual}.
    \item \textbf{Step 2.c (Interpretable Change Suggestion).} For each selected instance $x \in S^\star$, generate a domain-level edit description $\Delta(x_q, x)$ that specifies the transformations required to obtain $x$ from the query instance $x_q$. The edit must be interpretable, admissible under the constraints, and consistent with the cost function $c_\omega(x_q, x)$. 
\end{itemize}

Figure~\ref{fig:cf_two_stage_pipeline} illustrates the overall framework for counterfactual queries. In the offline stage, the system first performs \emph{1.a) Admissibility and cost parameterization} by examining the parameter space and selecting a configuration $\omega \in \mathcal{W}$, which instantiates the feasibility predicate $f_\omega$ and the change-cost function $c_\omega$. It then performs \emph{1.b) Constraint-aware index construction}, building a reusable index $\mathcal{I}$ over $\mathcal{D}$ to support efficient retrieval under these instance-level semantics. The index design is determined by the geometry induced by $\omega$. When $\omega$ restricts changes to a small mutable subspace or emphasizes a limited number of weighted dimensions, tree-based indices are typically effective. When $\omega$ allows greater flexibility or operates over high-dimensional representations, graph-based approximate nearest neighbor indices are more suitable.

In the online stage, given the prebuilt index $\mathcal{I}$, the system proceeds in three steps. First, \emph{2.a) Cost-aware candidate generation} retrieves and ranks feasible instances under $c_\omega$ and $f_\omega$, producing a candidate pool $\mathcal{A}_\omega(q)$ (green nodes) for query $q$ (orange node). Next, \emph{2.b) Set selection under $\Omega$} selects the final response set $S^\star \subseteq \mathcal{A}_\omega(q)$ (blue nodes) to optimize the specified set-level objective. Finally, \emph{2.c) Interpretable change selection} returns $S^\star$ together with domain-level edit descriptions $\Delta(q,x)$ for each selected instance $x \in S^\star$.